\structsection{ВВЕДЕНИЕ}

Утверждение «через точку вне прямой проходят по крайней мере две прямые, не пересекающие её» нельзя проверить на листе бумаги. Чертёж подчиняется евклидовой геометрии, и любая попытка нарисовать две непересекающиеся прямые через одну точку заканчивается тем, что они где-нибудь сходятся. Поэтому геометрию Лобачевского строят аксиоматически, а каждое доказательство опирается только на аксиомы и уже доказанные теоремы.

Модель Кэли~--- Клейна позволяет посмотреть на ту же геометрию с другой стороны. Точками в ней служат точки круга, прямыми~--- хорды, и все аксиомы абсолютной геометрии вместе с аксиомой Лобачевского в ней выполняются. Значит, каждая теорема геометрии Лобачевского верна и в модели, и её можно увидеть на конкретной фигуре. Обратное неверно: рисунок в модели ничего не доказывает, он лишь показывает один из случаев. В этой работе модель используется в обоих качествах. Как интерпретация она показывает, что означает утверждение. Как вычислительный инструмент она даёт явные формулы для угла параллельности, расстояния до прямой, углов четырёхугольника Саккери.

Компьютерная модель добавляет к этому движение. Если точку на рисунке можно передвинуть, становится видно, какие величины меняются вместе с ней, а какие остаются постоянными, и в какой момент пересекающиеся прямые становятся параллельными, а затем расходящимися. Чтобы числа на экране совпадали с теоремами, все построения в работе выведены из формул модели и проверены численно.

\textbf{Объект исследования}: планиметрия и простейшие объекты стереометрии Лобачевского.

\textbf{Предмет исследования}: интерпретация утверждений геометрии Лобачевского в модели Кэли~--- Клейна и их компьютерная визуализация.

\textbf{Цель работы}: построить в модели Кэли~--- Клейна интерпретации основных утверждений планиметрии Лобачевского, вывести для них явные формулы и создать по этим формулам интерактивные компьютерные модели.

Задачи работы:
\begin{enumerate}
  \item описать модель Кэли~--- Клейна на языке абсолюта: полюс и поляра, сложное отношение, π-движения, симметрия относительно прямой;
  \item интерпретировать в модели параллельные и расходящиеся прямые, угол параллельности, общий перпендикуляр и вывести формулу Лобачевского для угла параллельности;
  \item интерпретировать утверждения о треугольниках и четырёхугольниках, которые являются отрицаниями утверждений, равносильных V постулату;
  \item показать, как модель доказывает непротиворечивость геометрии Лобачевского относительно евклидовой и независимость аксиомы параллельности;
  \item распространить построение на пучки прямых, связки прямых и простейшие поверхности пространства Лобачевского;
  \item создать интерактивную веб-страницу с моделями плоскости и пространства Лобачевского;
  \item составить и решить задачи по теме работы.
\end{enumerate}

\textbf{Методы исследования}: аксиоматический метод, метод моделей, аналитический метод в однородных координатах, численная проверка формул.

Собственные результаты автора: вывод формул разделов 2 и~3 в модели Кэли~--- Клейна (угол параллельности с вершиной в центре круга, расстояние от точек прямой до другой прямой, углы четырёхугольника Саккери, граничная точка в отрицании предложения Лежандра); теорема о траекториях пучков и связок (теоремы~\ref{thm:traektoriya} и~\ref{thm:svyazka}) с доказательством, в котором три случая разбираются одним рассуждением; построение комнаты с полом-орисферой (теорема~\ref{thm:komnata}); веб-страница «Круг Кэли~--- Клейна» с моделями плоскости и пространства; решения задач раздела~7. Определения и классические теоремы приводятся по учебникам Л.\,С.~Атанасяна~[1], Л.\,С.~Атанасяна и В.\,Т.~Базылева~[2], Н.\,В.~Ефимова~[3].

Работа состоит из введения, семи разделов, заключения, списка использованных источников и приложения. Разделы~1--4 посвящены плоскости Лобачевского. Разделы~5 и~6 переносят те же методы на пространство: пучки и связки прямых, сферу, орисферу и эквидистантную поверхность, многогранники-«комнаты». Раздел~7 содержит задачи с решениями.
